%!TEX root = ../template.tex
%%%%%%%%%%%%%%%%%%%%%%%%%%%%%%%%%%%%%%%%%%%%%%%%%%%%%%%%%%%%%%%%%%%%
%% abstract-en.tex
%% NOVA thesis document file
%%
%% Abstract in English
%%%%%%%%%%%%%%%%%%%%%%%%%%%%%%%%%%%%%%%%%%%%%%%%%%%%%%%%%%%%%%%%%%%%

\typeout{NT FILE abstract-en.tex}%

Conflict-Free Replicated Data Types (CRDTs) are a key abstraction for modern distributed systems, allowing data to be updated concurrently across multiple replicas while ensuring Strong Eventual Consistency without centralized coordination. Despite these guarantees, the correct design of CRDTs remains challenging and is prone to errors that may compromise data integrity. Automated verification tools based on SMT solvers, support rapid prototyping and validation of convergence properties, but they face significant limitations when reasoning about recursive definitions, inductive properties, and complex global state invariants, often resulting in inconclusive verification outcomes.

This dissertation addresses these challenges by proposing an integrated framework for the modular and formal development of CRDTs that combines several deductive-based backend tools. The framework leverages the automation provided by VeriFx for checking algebraic convergence properties, while relying on the expressive power of the Why3 platform to ensure correctness guarantees through explicit contracts and invariants. A formal translation from VeriFx specifications to WhyML is introduced, enabling the use of deductive verification to enforce correct-by-construction semantics where automation alone is insufficient. Preliminary work, from translating and verifying fundamental CRDTs such as the Positive-Negative Counter in both VeriFx and Why3, indicate that specifications can be interchanged between tools without loss of correctness. The translation reveals the invariants and structural properties that need to be made explicit in a deductive setting, motivating the integration of multiple verification techniques.

% Abstract keywords
\keywords{
  CRDTs \and
  Formal Verification \and
  Why3 \and
  VeriFx
}
